\section{Conclusion}
\label{sec:experiments_conclusion}

The experiments reported in this chapter apply \gls{car} and its metrics to five Atari environments.
The surrogate policies reproduce the behavior of the initial policies with an average \BehaviorReconstructionName of $0.94 \pm 0.02$.
The latent clusterings are stable across independently trained surrogate policies, with a \ClusteringUniversalityName of about $0.5$ from three clusters onward, and they are largely independent of the observation and action spaces, with an average \AbstractionScoreName of $0.99 \pm 0.00$.
Together, these results provide empirical support for the \RepresentationConvergenceHypothesisName and, more broadly, for the hypothesis-driven approach underlying \gls{car}.

Several methodological choices have not yet been evaluated.
These include the \gls{nn} architecture, the layer selected for analysis, the distance used to compare embeddings, and the clustering algorithm.
Their impact should be studied jointly to determine how much the discovered \gls{cs} depends on the analysis procedure rather than on the policy itself.

The evaluation is also limited in scope.
It considers only five Atari games, and the semantic labels attached to the clusters in Figure~\ref{fig:semantic_granularity} were assigned by hand, which is time-consuming and subjective.
Extending the study to more diverse environments and developing human-centered evaluation protocols would strengthen the validity of the resulting explanations.
